\clearpage\SourceLeaf{529}\renewcommand{\BookRunningTitle}{Bemerkungen.}
\pdfbookmark[0]{Bemerkungen.}{bemerkungen-second-impression}
\begin{center}\textbf{BEMERKUNGEN.}\end{center}
\BookRule{.16\textwidth}
In diesen zweiten Band von \textsc{Gauss}’ Werken habe ich alle Abhandlungen, Aufsätze und Tafeln aus dem Gebiet der höheren Arithmetik aufgenommen, soweit sie nicht schon in den sieben Abschnitten der \textit{Disqu. Arithm.} enthalten sind: die fünf in den ‘\textit{Commentationes societatis regiae scientiarum Gottingensis}’ (in Quart) veröffentlichten Abhandlungen, die in den ‘\textit{Göttingischen Gelehrten Anzeigen}’ (in Oktav) erschienenen Besprechungen sowohl dieser eigenen als auch einiger fremder Schriften (von \textsc{Gauss} nicht unterzeichnet, aber durch die Akten der Göttinger Universitätsbibliothek als seine Arbeiten bestätigt), sowie eine Auswahl aus dem handschriftlichen Nachlass.
\par Beim zweiten Abdruck habe ich noch die Tafeln ‘Circuli quadratura nova’, ‘Zur Berechnung der Logarithmen’, ‘Quadratorum myrias prima’, ‘Indizes der Primzahlen im höheren Zahlenbereich’, ‘Hilfstafel zur Lösung der unbestimmten Gleichung $A=fxx+gyy$ mit dem Ausschlussverfahren’, ferner ‘Sectio octava’, soweit sie niedergeschrieben ist, und schließlich zwei Briefe von \textsc{Gauss} an \textsc{Dirichlet} als wesentliche Stücke der Geschichte der höheren Arithmetik hinzugefügt.
\par Damit die Gegenstände in einem so umfangreichen Band leichter zu überblicken sind, werden die Lehrsätze einheitlich durch den Druck hervorgehoben. Um den Gebrauch sowohl der älteren Ausgaben als auch der vorliegenden zu erleichtern, sind bei Verweisen auf die Disq. Arithm. die Seitenzahlen durch die Artikelnummern ersetzt, bei Angaben von Abhandlungen die Veröffentlichungsorte durch die eigenen Titel. Die Anmerkung, die ursprünglich Artikel 2 der Abhandlung ‘\textit{Theorematis arithmetici demonstratio nova}’
